\chapter{\ploren{Wykresy i prezentacja danych}{Plots and data presentation}}
\label{app:wykresy}

\begin{quote}
\itshape
\ploren
  {Ten załącznik zawiera przykłady instruktażowe przeznaczone do usunięcia lub zastąpienia wynikami autora. Dane są demonstracyjne i nie stanowią rezultatów rzeczywistych badań.}
  {This appendix contains instructional examples that should be removed or replaced with the author's own results. The data are illustrative and do not constitute results of an actual study.}
\end{quote}

Wykres powinien odpowiadać na określone pytanie i umożliwiać poprawne odczytanie danych. Nie należy traktować go wyłącznie jako elementu dekoracyjnego. Każdy wykres wymaga omówienia w tekście, podpisu, etykiety oraz jednoznacznych nazw osi i jednostek. Zasady dotyczące formatów, czytelności, źródeł i plików edytowalnych przedstawiono w Załączniku~\ref{app:rysunki}; obowiązują one również w tej części.

\section{\ploren{Formaty przechowywania danych}{Data storage formats}}

CSV jest najczęściej stosowanym formatem wymiany niewielkich i średnich tabel danych. Jest tekstowy, łatwy do kontroli wersji i obsługiwany przez MATLAB, Python, arkusze kalkulacyjne oraz \texttt{pgfplots}. Należy jednak jawnie ustalić separator kolumn, znak dziesiętny, kodowanie i obecność nagłówka. W pliku przeznaczonym do automatycznego przetwarzania bezpiecznym wyborem są przecinek jako separator kolumn oraz kropka jako znak dziesiętny.

W praktyce spotyka się również:

\begin{itemize}
  \item TSV — tekstowe kolumny rozdzielone tabulatorem; separator tabulatora wskazuje się podczas wczytywania danych,
  \item DAT lub TXT — dane rozdzielone spacjami albo zapisane w formacie właściwym dla danego programu; dla białych znaków można użyć \verb|col sep=space|,
  \item XLSX lub ODS — wygodne do ręcznego przeglądania, ale mniej odpowiednie do kontroli wersji; dane do wykresu warto wyeksportować do CSV,
  \item MAT, NPZ, HDF5 lub Parquet — przydatne dla dużych tablic, wielu kanałów i metadanych; przed użyciem w LaTeX zwykle przygotowuje się skrypt tworzący CSV albo gotowy wykres,
  \item JSON — właściwy dla danych hierarchicznych i wyników aplikacji, lecz przed narysowaniem wykresu zazwyczaj wymaga przekształcenia do tabeli.
\end{itemize}

Nie należy ręcznie kopiować długich kolumn do kodu LaTeX. Plik źródłowy powinien zachować dane z możliwie pełną dokładnością, natomiast sposób zaokrąglania wartości należy określić dopiero podczas tworzenia tabeli lub wykresu. Komentarze i metadane muszą być zapisane tak, aby nie zostały omyłkowo odczytane jako obserwacje.

Przykładowy plik tekstowy może zostać wczytany następująco:

\begin{lstlisting}[numbers=none]
\addplot table[
  x=frequency_Hz,
  y=gain_dB,
  col sep=space,
  comment chars={\#}
]{data/charakterystyka.dat};
\end{lstlisting}

\section{\ploren{Wykres tworzony bezpośrednio w LaTeX}{Plot generated directly in LaTeX}}

Pakiet \texttt{pgfplots} umożliwia obliczenie i narysowanie funkcji bez używania zewnętrznego programu. Rysunek~\ref{fig:wykres-funkcji} przedstawia prosty przebieg sinusoidalny. Taki sposób jest przydatny dla zależności analitycznych i niewielkich przykładów, ponieważ tekst oraz symbole zachowują styl dokumentu.

\begin{figure}[htbp]
  \centering
  \begin{tikzpicture}
    \begin{axis}[
      width=0.86\textwidth,
      height=0.46\textwidth,
      xlabel={Czas $t$ (\si{\second})},
      ylabel={Amplituda (y)},
      xmin=0,xmax=6.283,
      ymin=-1.15,ymax=1.15,
      grid=both,
      minor tick num=1,
      domain=0:6.283,
      samples=160
    ]
      \addplot[thesislinkblue,very thick] {sin(deg(x))};
    \end{axis}
  \end{tikzpicture}
  \caption{\ploren{Przykładowy przebieg funkcji sinusoidalnej}{Example sinusoidal waveform}}
  \label{fig:wykres-funkcji}
\end{figure}

\begin{lstlisting}[caption={\ploren{Wykres funkcji w pakiecie pgfplots}{Function plot using pgfplots}}]
\begin{tikzpicture}
  \begin{axis}[
    xlabel={Czas \(t\) (\si{\second})},
    ylabel={Amplituda \(y\)},
    grid=both,
    domain=0:6.283,
    samples=160
  ]
    \addplot[thesislinkblue,very thick]
      {sin(deg(x))};
  \end{axis}
\end{tikzpicture}
\end{lstlisting}

Parametr \texttt{domain} wyznacza zakres argumentu, natomiast \texttt{samples} liczbę punktów użytych do narysowania krzywej. Zbyt mała liczba próbek może powodować widoczne załamania, a nadmierna niepotrzebnie wydłuża kompilację.

\section{\ploren{Wykres z danych zapisanych w CSV}{Plot from CSV data}}

Dane pomiarowe lub wyniki obliczeń najlepiej przechowywać poza plikiem \texttt{.tex}. Rysunek~\ref{fig:wykres-csv} jest generowany z pliku \texttt{data/odpowiedz.csv}. Zmiana danych nie wymaga wtedy ręcznego poprawiania współrzędnych w kodzie wykresu.

\begin{figure}[htbp]
  \centering
  \begin{tikzpicture}
    \begin{axis}[
      width=0.88\textwidth,
      height=0.50\textwidth,
      xlabel={Czas $t$ (\si{\second})},
      ylabel={Wartość znormalizowana},
      xmin=0,xmax=2,
      ymin=0,ymax=1.08,
      grid=major,
      legend pos=south east,
      legend cell align=left
    ]
      \addplot[thesislinkblue,very thick,mark=*]
        table[x=time,y=reference,col sep=comma]
        {chapters/11-zalacznik-d-wykresy/data/odpowiedz.csv};
      \addlegendentry{Wartość zadana}
      \addplot[thesiscodestring,very thick,dashed,mark=square*]
        table[x=time,y=measured,col sep=comma]
        {chapters/11-zalacznik-d-wykresy/data/odpowiedz.csv};
      \addlegendentry{Wartość zmierzona}
    \end{axis}
  \end{tikzpicture}
  \caption{\ploren{Porównanie wartości zadanej i zmierzonej na podstawie danych CSV}{Comparison of reference and measured values from CSV data}}
  \label{fig:wykres-csv}
\end{figure}

\begin{lstlisting}[caption={\ploren{Wczytywanie dwóch serii z pliku CSV}{Loading two series from a CSV file}}]
\addplot[blue,thick,mark=*]
  table[x=time,y=reference,col sep=comma]
  {data/odpowiedz.csv};
\addlegendentry{Wartość zadana}

\addplot[red,thick,dashed,mark=square*]
  table[x=time,y=measured,col sep=comma]
  {data/odpowiedz.csv};
\addlegendentry{Wartość zmierzona}
\end{lstlisting}

Serie rozróżniono jednocześnie kolorem, rodzajem linii i znacznikiem. Dzięki temu wykres pozostaje czytelny po wydruku w skali szarości. Plik CSV powinien zawierać nagłówki opisujące kolumny, a separator określa się opcją \texttt{col sep}.

\section{\ploren{Aproksymacja i interpolacja danych}{Data approximation and interpolation}}

\subsection{\ploren{Różnica między aproksymacją i interpolacją}{Difference between approximation and interpolation}}

Aproksymacja wyznacza model, który opisuje ogólną zależność w danych, ale nie musi przechodzić przez każdy punkt pomiarowy. Stosuje się ją między innymi do estymacji parametrów, ograniczenia wpływu szumu i opisu trendu. Interpolacja wyznacza wartości pomiędzy znanymi punktami i w typowym przypadku przechodzi przez wszystkie dane wejściowe. Nie powinna być mylona z ekstrapolacją poza zakres pomiarów.

Na rysunku~\ref{fig:aproksymacja-interpolacja} przedstawiono oba przypadki obok siebie. Punkty oznaczają dane pomiarowe. Na rysunku~\ref{fig:aproksymacja} linia ciągła przedstawia liniowy model aproksymacyjny (y=\num{2.049}x+\num{0.895}), którego współczynniki wyznaczono wcześniej w skrypcie obliczeniowym. Na rysunku~\ref{fig:interpolacja} linia przerywana łączy wartości uzyskane przez interpolację liniową wewnątrz zakresu danych.

\begin{figure}[htbp]
  \centering
  \begin{subfigure}[t]{0.48\textwidth}
    \centering
    \begin{tikzpicture}
      \begin{axis}[
        width=\linewidth,
        height=0.78\linewidth,
        xlabel={Wielkość wejściowa (x)},
        ylabel={Wielkość wyjściowa (y)},
        xmin=-0.2,xmax=5.2,
        ymin=0,ymax=12,
        grid=major,
        legend style={font=\footnotesize},
        legend pos=north west
      ]
        \addplot[only marks,mark=*,black,mark size=2.2pt]
          table[x=x,y=measurement,col sep=comma]
          {chapters/11-zalacznik-d-wykresy/data/aproksymacja.csv};
        \addlegendentry{Pomiary}
        \addplot[thesislinkblue,line width=1.2pt]
          table[x=x,y=fit,col sep=comma]
          {chapters/11-zalacznik-d-wykresy/data/aproksymacja.csv};
        \addlegendentry{Aproksymacja}
      \end{axis}
    \end{tikzpicture}
    \caption{Aproksymacja liniowa}
    \label{fig:aproksymacja}
  \end{subfigure}
  \hfill
  \begin{subfigure}[t]{0.48\textwidth}
    \centering
    \begin{tikzpicture}
      \begin{axis}[
        width=\linewidth,
        height=0.78\linewidth,
        xlabel={Wielkość wejściowa (x)},
        ylabel={Wielkość wyjściowa (y)},
        xmin=-0.2,xmax=5.2,
        ymin=0,ymax=12,
        grid=major,
        legend style={font=\footnotesize},
        legend pos=north west
      ]
        \addplot[thesiscodestring,dashed,line width=1.2pt]
          table[x=x,y=interpolated,col sep=comma]
          {chapters/11-zalacznik-d-wykresy/data/interpolacja.csv};
        \addlegendentry{Interpolacja}
        \addplot[only marks,mark=square*,black,mark size=2.2pt]
          table[x=x,y=measurement,col sep=comma]
          {chapters/11-zalacznik-d-wykresy/data/aproksymacja.csv};
        \addlegendentry{Pomiary}
      \end{axis}
    \end{tikzpicture}
    \caption{Interpolacja liniowa}
    \label{fig:interpolacja}
  \end{subfigure}
  \caption{Porównanie aproksymacji i interpolacji demonstracyjnych danych pomiarowych}
  \label{fig:aproksymacja-interpolacja}
\end{figure}

\begin{lstlisting}[caption={\ploren{Aproksymacja i interpolacja z danych CSV}{Approximation and interpolation from CSV data}}]
% Punkty pomiarowe
\addplot[only marks,mark=*,black]
  table[x=x,y=measurement,col sep=comma]
  {data/aproksymacja.csv};

% Wartości modelu aproksymacyjnego
\addplot[blue,line width=1.2pt]
  table[x=x,y=fit,col sep=comma]
  {data/aproksymacja.csv};

% Wartości obliczone przez interpolację
\addplot[red,dashed,line width=1.2pt]
  table[x=x,y=interpolated,col sep=comma]
  {data/interpolacja.csv};
\end{lstlisting}

Metodę wyznaczenia współczynników aproksymacji i wartości interpolowanych należy opisać w tekście oraz zachować w skrypcie MATLAB-a, Pythona lub innego środowiska. Współczynniki modelu, stopień wielomianu, rodzaj funkcji bazowej i sposób ważenia punktów nie mogą wynikać wyłącznie z ręcznego dopasowania wyglądu krzywej. Należy również sprawdzić, czy wybrany model ma uzasadnienie techniczne i czy nie powoduje nadmiernego dopasowania.

\subsection{\ploren{Ocena jakości aproksymacji}{Evaluation of approximation quality}}

Ocena modelu nie powinna opierać się wyłącznie na wyglądzie dopasowanej linii. Dla przedstawionych danych błąd średniokwadratowy wynosi $\mathrm{RMSE}=\num{0.155}$, a współczynnik determinacji $R^2=\num{0.998}$. Wartości te trzeba interpretować łącznie z jednostkami, zakresem danych, liczbą pomiarów i założeniami modelu.

Na rysunku~\ref{fig:reszty-aproksymacji} pokazano reszty, czyli różnice pomiędzy pomiarem i wartością modelu. Ich analiza pozwala zauważyć systematyczną zależność, zmienną wariancję albo obserwacje odstające, których sama wartość \(R^2\) może nie ujawnić.

\begin{figure}[htbp]
  \centering
  \begin{tikzpicture}
    \begin{axis}[
      width=0.82\textwidth,
      height=0.42\textwidth,
      xlabel={Wielkość wejściowa (x)},
      ylabel={Reszta $e=y-\hat{y}$},
      xmin=-0.2,xmax=5.2,
      ymin=-0.30,ymax=0.30,
      grid=major
    ]
      \addplot[black,thin] coordinates {(0,0) (5,0)};
      \addplot[only marks,mark=*,thesislinkblue,mark size=2.5pt]
        table[x=x,y=residual,col sep=comma]
        {chapters/11-zalacznik-d-wykresy/data/aproksymacja.csv};
    \end{axis}
  \end{tikzpicture}
  \caption{Reszty liniowego modelu aproksymacyjnego}
  \label{fig:reszty-aproksymacji}
\end{figure}

\begin{lstlisting}[caption={\ploren{Wykres reszt modelu}{Residual plot}}]
\addplot[black,thin]
  coordinates {(0,0) (5,0)};
\addplot[only marks,mark=*,blue]
  table[x=x,y=residual,col sep=comma]
  {data/aproksymacja.csv};
\end{lstlisting}

Wartość \(R^2\) bliska jedności nie dowodzi automatycznie poprawności modelu ani związku przyczynowego. Dla modeli nieliniowych i danych walidacyjnych mogą być potrzebne inne miary. Zwiększanie stopnia wielomianu tylko w celu poprawienia wskaźnika może prowadzić do nadmiernego dopasowania i niestabilnej ekstrapolacji.

\section{\ploren{Kilka wykresów w jednym układzie}{Multiple plots in one layout}}

Przykład z rysunku~\ref{fig:aproksymacja-interpolacja} wykorzystuje dwa środowiska \texttt{subfigure}, z których każde zajmuje \SI{48}{\percent} szerokości tekstu. Dzięki temu wykresy znajdują się obok siebie i mają wspólny podpis. Taki układ jest właściwy tylko wtedy, gdy opisy osi i legendy pozostają czytelne.

\begin{lstlisting}[caption={\ploren{Dwa wykresy obok siebie}{Two plots side by side}}]
\begin{figure}[htbp]
  \centering
  \begin{subfigure}[t]{0.48\textwidth}
    \centering
    \begin{tikzpicture}
      \begin{axis}[width=\linewidth]
        % pierwszy wykres
      \end{axis}
    \end{tikzpicture}
    \caption{Pierwszy wykres}
  \end{subfigure}
  \hfill
  \begin{subfigure}[t]{0.48\textwidth}
    \centering
    \begin{tikzpicture}
      \begin{axis}[width=\linewidth]
        % drugi wykres
      \end{axis}
    \end{tikzpicture}
    \caption{Drugi wykres}
  \end{subfigure}
  \caption{Wspólny podpis}
\end{figure}
\end{lstlisting}

Jeżeli wykresy są zbyt małe, należy umieścić je jeden pod drugim. Każdy element może wtedy wykorzystać całą szerokość tekstu:

\begin{lstlisting}[caption={\ploren{Dwa wykresy jeden pod drugim}{Two plots stacked vertically}}]
\begin{figure}[htbp]
  \centering
  \begin{subfigure}{\textwidth}
    \centering
    \begin{tikzpicture}
      \begin{axis}[width=0.85\textwidth]
        % pierwszy wykres
      \end{axis}
    \end{tikzpicture}
    \caption{Pierwszy wykres}
  \end{subfigure}

  \medskip

  \begin{subfigure}{\textwidth}
    \centering
    \begin{tikzpicture}
      \begin{axis}[width=0.85\textwidth]
        % drugi wykres
      \end{axis}
    \end{tikzpicture}
    \caption{Drugi wykres}
  \end{subfigure}
  \caption{Wspólny podpis}
\end{figure}
\end{lstlisting}

Można także zastosować kilka środowisk \texttt{axis} bez \texttt{subfigure}, pakietowe układy grupowe albo osobne rysunki. Pokazane rozwiązania nie są jedynymi możliwościami; wybór powinien zależeć od liczby wykresów, potrzeby wspólnego podpisu i czytelności w docelowym rozmiarze.

\section{\ploren{Grubość linii, znaczniki i style}{Line width, markers and styles}}

Styl serii określa się w opcjach polecenia \verb|\addplot|. Dostępne są między innymi predefiniowane grubości \texttt{thin}, \texttt{semithick}, \texttt{thick} i \texttt{very thick}. Dokładną wartość można podać przez \verb|line width=1.2pt|. Rodzaj linii określają opcje \texttt{solid}, \texttt{dashed}, \texttt{dotted} i \texttt{dashdotted}.

\begin{lstlisting}[caption={\ploren{Przykłady zmiany stylu serii danych}{Examples of changing a data-series style}}]
\addplot[thin,solid,blue] ...;
\addplot[very thick,dashed,red] ...;
\addplot[line width=1.2pt,dashdotted,black] ...;
\addplot[
  semithick,
  mark=square*,
  mark size=2.5pt,
  mark options={fill=white},
  color=blue
] ...;
\end{lstlisting}

Można również ustawić \texttt{opacity}, własny wzór kreskowania, kolor znacznika, rozmiar symboli i częstotliwość ich wyświetlania za pomocą \texttt{mark repeat}. Styl powinien być spójny w całej pracy. Warto zdefiniować powtarzające się zestawy opcji za pomocą \verb|\pgfplotsset|, zamiast kopiować je do każdego wykresu:

\begin{lstlisting}[caption={\ploren{Definicja własnych stylów wykresu}{Definition of custom plot styles}}]
\pgfplotsset{
  pomiar/.style={
    only marks,
    mark=*,
    mark size=2.2pt,
    color=black
  },
  model/.style={
    line width=1.2pt,
    color=blue
  }
}

\addplot[pomiar] ...;
\addplot[model] ...;
\end{lstlisting}

Kolory można definiować poleceniem \verb|\definecolor|, ale nie należy opierać rozróżnienia serii wyłącznie na barwie. Grubsza linia nie powinna służyć do sugerowania większej ważności wyniku, jeśli taka hierarchia nie została wyjaśniona.

\section{\ploren{Niepewność i rozrzut danych}{Uncertainty and data variability}}

Pojedyncza linia nie informuje o niepewności ani zmienności wyniku. Jeżeli są one istotne dla wniosków, należy zastosować słupki błędów, pasmo ufności, wykres pudełkowy albo przedstawić pojedyncze obserwacje. Na rysunku~\ref{fig:wykres-niepewnosc} słupki pokazują demonstracyjną niepewność wartości zmierzonej.

\begin{figure}[htbp]
  \centering
  \begin{tikzpicture}
    \begin{axis}[
      width=0.86\textwidth,
      height=0.48\textwidth,
      xlabel={Czas $t$ (\si{\second})},
      ylabel={Wartość zmierzona},
      xmin=0,xmax=2,
      ymin=0,ymax=1.08,
      grid=major
    ]
      \addplot[only marks,mark=*,thesislinkblue,
        error bars/.cd,y dir=both,y explicit]
        table[x=time,y=measured,y error=uncertainty,col sep=comma]
        {chapters/11-zalacznik-d-wykresy/data/odpowiedz.csv};
    \end{axis}
  \end{tikzpicture}
  \caption{\ploren{Przykład przedstawienia wartości wraz z niepewnością}{Example presentation of values with uncertainty}}
  \label{fig:wykres-niepewnosc}
\end{figure}

\begin{lstlisting}[caption={\ploren{Słupki błędów z wartości zapisanych w CSV}{Error bars from values stored in CSV}}]
\addplot[
  only marks,
  mark=*,
  error bars/.cd,
  y dir=both,
  y explicit
]
table[
  x=time,
  y=measured,
  y error=uncertainty,
  col sep=comma
]{data/odpowiedz.csv};
\end{lstlisting}

W podpisie lub tekście trzeba wyjaśnić, co oznaczają słupki błędów: odchylenie standardowe, błąd standardowy, przedział ufności czy niepewność pomiaru. Samo narysowanie słupków bez tej informacji jest niewystarczające.

\section{\ploren{Wykres słupkowy}{Bar chart}}

Wykres słupkowy jest odpowiedni do porównania odrębnych kategorii, ale nie powinien zastępować wykresu przebiegu ciągłego. Rysunek~\ref{fig:wykres-slupkowy} przedstawia przykład porównania czasu obliczeń trzech wariantów. Wartości umieszczono nad słupkami, dzięki czemu wykres pozostaje jednoznaczny bez odczytywania ich wyłącznie z wysokości.

\begin{figure}[htbp]
  \centering
  \begin{tikzpicture}
    \begin{axis}[
      ybar,
      width=0.72\textwidth,
      height=0.48\textwidth,
      ylabel={Czas obliczeń (\si{\milli\second})},
      symbolic x coords={A,B,C},
      xtick=data,
      ymin=0,ymax=50,
      nodes near coords,
      bar width=22pt,
      grid=major
    ]
      \addplot[fill=thesislinkblue!55,draw=black]
        coordinates {(A,42) (B,31) (C,24)};
    \end{axis}
  \end{tikzpicture}
  \caption{\ploren{Przykładowe porównanie czasu obliczeń trzech wariantów}{Example comparison of computation time for three variants}}
  \label{fig:wykres-slupkowy}
\end{figure}

Oś wykresu słupkowego powinna zwykle rozpoczynać się od zera, ponieważ długość słupka reprezentuje wartość. Skrócenie osi może wizualnie wyolbrzymić niewielkie różnice.

\section{\ploren{Skala logarytmiczna i charakterystyka częstotliwościowa}{Logarithmic scale and frequency response}}

Skalę logarytmiczną stosuje się, gdy wartości obejmują wiele rzędów wielkości albo gdy wynika to z przyjętej metody analizy. Jest ona typowa dla charakterystyk częstotliwościowych, wykresów Bodego, widm oraz zależności potęgowych. Rysunek~\ref{fig:skala-logarytmiczna} przedstawia demonstracyjną charakterystykę amplitudową układu dolnoprzepustowego o częstotliwości granicznej \(f_c=\SI{100}{\hertz}\).

\begin{figure}[htbp]
  \centering
  \begin{tikzpicture}
    \begin{semilogxaxis}[
      width=0.86\textwidth,
      height=0.48\textwidth,
      xlabel={Częstotliwość $f$ (\si{\hertz})},
      ylabel={Wzmocnienie $K$ (\si{\decibel})},
      xmin=10,xmax=10000,
      ymin=-42,ymax=2,
      grid=both,
      domain=10:10000,
      samples=180
    ]
      \addplot[thesislinkblue,line width=1.2pt]
        {-10*ln(1+(x/100)^2)/ln(10)};
      \addplot[black,dashed,thin]
        coordinates {(100,-42) (100,2)};
    \end{semilogxaxis}
  \end{tikzpicture}
  \caption{Przykład charakterystyki amplitudowej ze skalą logarytmiczną częstotliwości}
  \label{fig:skala-logarytmiczna}
\end{figure}

\begin{lstlisting}[caption={\ploren{Wykres z logarytmiczną osią częstotliwości}{Plot with a logarithmic frequency axis}}]
\begin{semilogxaxis}[
  xlabel={Częstotliwość \(f\) (\si{\hertz})},
  ylabel={Wzmocnienie \(K\) (\si{\decibel})},
  xmin=10,
  xmax=10000,
  grid=both
]
  \addplot[blue,line width=1.2pt]
    {-10*ln(1+(x/100)^2)/ln(10)};
\end{semilogxaxis}
\end{lstlisting}

Pakiet udostępnia również środowiska \texttt{semilogyaxis} i \texttt{loglogaxis}. Nie wolno stosować skali logarytmicznej dla zera ani wartości ujemnych. Rodzaj skali musi być widoczny na osi i uwzględniony podczas interpretacji nachylenia oraz odległości pomiędzy punktami.

\section{\ploren{MATLAB, Python i inne narzędzia}{MATLAB, Python and other tools}}

Wykresy przygotowane w MATLAB-ie, Pythonie lub innym środowisku można wykorzystać na dwa profesjonalne sposoby. Pierwszy polega na zapisaniu danych do CSV i utworzeniu wykresu w \texttt{pgfplots}. Zapewnia to spójne czcionki i styl, ale może wydłużyć kompilację przy bardzo dużych zbiorach danych.

Drugi sposób polega na przygotowaniu kompletnego wykresu w programie źródłowym i eksporcie do PDF. Jest zalecany dla złożonych wykresów, dużej liczby punktów oraz typów grafiki, których \texttt{pgfplots} nie obsługuje wygodnie. PNG należy pozostawić głównie dla wyników rastrowych, takich jak mapy cieplne, obrazy, spektrogramy lub wizualizacje macierzy. Nie należy wykonywać zrzutu ekranu, jeśli program umożliwia prawidłowy eksport.

Skrypt generujący wykres, dane i informacje o użytych wersjach bibliotek powinny pozostać w katalogu zasobów rozdziału. Wykres musi dać się odtworzyć bez ręcznej zmiany wartości, kolorów lub podpisów w programie graficznym.

Poniższe skrypty wczytują ten sam plik CSV, wyznaczają liniowy model aproksymacyjny, obliczają RMSE i \(R^2\), a następnie eksportują wykres wektorowy PDF. Są to kompletne pliki znajdujące się w katalogu \texttt{scripts}.

Skrypt Pythona wymaga bibliotek \texttt{numpy} i \texttt{matplotlib}, wymienionych w pliku \texttt{scripts/requirements.txt}. Można je zainstalować w używanym środowisku wirtualnym poleceniem \verb|python -m pip install -r scripts/requirements.txt|. Zależności projektu powinny być instalowane w środowisku przeznaczonym dla danej pracy, a nie dopisywane ręcznie na komputerze osoby kompilującej dokument.

\lstinputlisting[
  style=thesis-python,
  caption={\ploren{Aproksymacja i eksport wykresu w Pythonie}{Approximation and plot export in Python}}
]{chapters/11-zalacznik-d-wykresy/scripts/wykres.py}

\lstinputlisting[
  style=thesis-matlab,
  caption={\ploren{Aproksymacja i eksport wykresu w MATLAB-ie}{Approximation and plot export in MATLAB}}
]{chapters/11-zalacznik-d-wykresy/scripts/wykres.m}

Wersje języka i bibliotek należy zapisać w dokumentacji projektu albo pliku środowiska. Sam plik PDF nie wystarcza do odtworzenia wyniku, jeżeli nie zachowano danych, kodu i konfiguracji.

\section{\ploren{Zaawansowane sposoby prezentacji danych}{Advanced data presentation}}

\subsection{\ploren{Histogram}{Histogram}}

Histogram przedstawia przybliżony rozkład obserwacji przez zliczenie wartości w przedziałach. Liczba i szerokość przedziałów wpływają na wygląd, dlatego nie należy dobierać ich wyłącznie w celu uzyskania oczekiwanego obrazu. Rysunek~\ref{fig:histogram} wykorzystuje pojedynczą kolumnę pliku CSV.

\begin{figure}[htbp]
  \centering
  \begin{tikzpicture}
    \begin{axis}[
      ybar interval,
      width=0.78\textwidth,
      height=0.46\textwidth,
      xlabel={Wartość zmierzona},
      ylabel={Liczba obserwacji},
      grid=major
    ]
      \addplot[hist={bins=8},fill=thesislinkblue!55,draw=black]
        table[y=value,col sep=comma]
        {chapters/11-zalacznik-d-wykresy/data/proby.csv};
    \end{axis}
  \end{tikzpicture}
  \caption{Histogram demonstracyjnych wyników pomiarów}
  \label{fig:histogram}
\end{figure}

\begin{lstlisting}[caption={\ploren{Histogram danych z pliku CSV}{Histogram of data from a CSV file}}]
\addplot[
  hist={bins=8},
  fill=blue!55,
  draw=black
]
table[y=value,col sep=comma]
{data/proby.csv};
\end{lstlisting}

\subsection{\ploren{Mapa cieplna}{Heat map}}

Mapa cieplna przedstawia wartość za pomocą koloru w dwóch wymiarach przestrzennych lub parametrycznych. Musi zawierać skalę barwną z nazwą wielkości i jednostką. Rysunek~\ref{fig:mapa-cieplna} pokazuje demonstracyjny rozkład temperatury.

\begin{figure}[htbp]
  \centering
  \begin{tikzpicture}
    \begin{axis}[
      width=0.72\textwidth,
      height=0.58\textwidth,
      xlabel={Położenie $x$ (\si{\centi\metre})},
      ylabel={Położenie $y$ (\si{\centi\metre})},
      view={0}{90},
      colorbar,
      colorbar style={ylabel={Temperatura (\si{\degreeCelsius})}},
      colormap/hot2
    ]
      \addplot[matrix plot*,point meta=explicit]
        table[x=x,y=y,meta=value,col sep=comma]
        {chapters/11-zalacznik-d-wykresy/data/mapa-cieplna.csv};
    \end{axis}
  \end{tikzpicture}
  \caption{Przykładowa mapa cieplna rozkładu temperatury}
  \label{fig:mapa-cieplna}
\end{figure}

\begin{lstlisting}[caption={\ploren{Mapa cieplna z wartościami jawnymi}{Heat map with explicit values}}]
\begin{axis}[
  view={0}{90},
  colorbar,
  colormap/hot2
]
  \addplot[matrix plot*,point meta=explicit]
    table[x=x,y=y,meta=value,col sep=comma]
    {data/mapa-cieplna.csv};
\end{axis}
\end{lstlisting}

Paleta barw powinna być czytelna dla osób z zaburzeniami rozpoznawania kolorów i po wydruku w skali szarości. Palet tęczowych należy unikać, ponieważ mogą tworzyć sztuczne granice i utrudniać ocenę zmian wartości.

\subsection{\ploren{Spektrogram}{Spectrogram}}

Spektrogram jest mapą wartości w dziedzinie czasu i częstotliwości. Oprócz osi i skali barw trzeba podać sposób wyznaczenia widma: długość okna, rodzaj okna, nakładanie segmentów, częstotliwość próbkowania i sposób skalowania amplitudy. Rysunek~\ref{fig:spektrogram} jest niewielkim przykładem demonstracyjnym.

\begin{figure}[htbp]
  \centering
  \begin{tikzpicture}
    \begin{axis}[
      width=0.78\textwidth,
      height=0.52\textwidth,
      xlabel={Czas $t$ (\si{\second})},
      ylabel={Częstotliwość $f$ (\si{\hertz})},
      view={0}{90},
      colorbar,
      colorbar style={ylabel={Moc znormalizowana}},
      colormap/hot2
    ]
      \addplot[matrix plot*,point meta=explicit]
        table[x=time,y=frequency,meta=power,col sep=comma]
        {chapters/11-zalacznik-d-wykresy/data/spektrogram.csv};
    \end{axis}
  \end{tikzpicture}
  \caption{Demonstracyjny spektrogram sygnału}
  \label{fig:spektrogram}
\end{figure}

\begin{lstlisting}[caption={\ploren{Spektrogram zapisany jako tabela wartości}{Spectrogram stored as a table of values}}]
\addplot[matrix plot*,point meta=explicit]
  table[
    x=time,
    y=frequency,
    meta=power,
    col sep=comma
  ]
  {data/spektrogram.csv};
\end{lstlisting}

Dla dużych macierzy spektrogram lub mapę cieplną zwykle lepiej wygenerować w MATLAB-ie albo Pythonie i wyeksportować do PNG lub PDF. W pracy należy zachować dane oraz parametry użyte do utworzenia obrazu.

\section{\ploren{Dobór rodzaju wykresu i typowe błędy}{Choosing a plot type and common mistakes}}

Wykres liniowy służy do prezentowania przebiegu lub uporządkowanej zależności, punktowy do pokazania obserwacji i korelacji, słupkowy do porównania kategorii, a histogram do przedstawienia rozkładu. Skala logarytmiczna jest właściwa przy danych obejmujących wiele rzędów wielkości, ale musi zostać jasno oznaczona.

Należy unikać wykresów trójwymiarowych dla danych, które można czytelniej przedstawić w dwóch wymiarach, podwójnych osi pionowych bez mocnego uzasadnienia, nadmiernie wygładzonych przebiegów, niewyjaśnionych skrótów, zbyt małej czcionki i ręcznie dobranych granic osi sugerujących nieistniejące różnice. Wykres powinien prezentować dane uczciwie i umożliwiać ich ocenę, a nie jedynie wspierać oczekiwaną tezę.
